\scanpage{267}
\mspage{8}

Une façon plus claire de voir ceci est la suivante. Soit $J\deleted{\unsure{*}}$ une catégorie \inserted{intervalle, i.e.} isomorphe à un intervalle \inserted{standard}. Alors il y a exactement \inserted{(si $I\ne e$)} \uline{deux} couples $(I,\Theta)$ d'un intervalle standard $I$ et d'un isom. $\Theta:I\xrightarrow{\sim}J$.

Ces deux isomorphismes induisent sur $J$ deux ordres \deleted{\unsure{tot}} totaux \deleted{inverses} \inserted{opposés} l'un de l'autre, \deleted{et sont} Chaque couple est déterminé de façon unique par ledit ordre \deleted{total}, qu'on appellera une \uline{orientation} de l'intervalle $J$.

% slanted
\margin{\uline{NB} On parlera de l'intervalle \uline{orienté} \uline{opposé} $I^{\mathrm{opp}}$ d'un intervalle \uline{orienté} $I$, qu'il faut se garder de confondre avec la catégorie opposée $I^{\unsure{\circ|\mathrm{o}}}$. $I^{\mathrm{opp}}$ a $\mathord{=}$ catégorie \unsure{ss-jacente} que $I$, mais avec une orientation opposée.}

Ainsi $J$ a deux orientations, et \deleted{pour} chacune détermine un \uline{isomorphisme} \uline{orienté} (i.e. respectant les orientations) de $J$ sur un intervalle standard. \unsure{Et|Si} les deux intervalles sont les mêmes, ssi. il existe un automorphisme de $J$ \inserted{(nécessairement unique)} qui échange ses deux orientations, ou ce qui revient au $\mathord{=}$, qui échange entre elles les deux extrémités. (Ainsi les catégories opposées l'une de l'autre

\diagram{\begin{tikzcd}[ampersand replacement=\&,column sep=small,row sep=small]
 \& \cdot\arrow[dl]\arrow[dr] \& \& \& \cdot \& \\
 \cdot \& \& \cdot \& \cdot\arrow[ur] \& \& \cdot\arrow[ul] \& {,}
\end{tikzcd}}

sont chacune des intervalles symétriques.) \inserted{Choisir une orientation revient \unsure{ainsi|aussi} à choisir l'origine parmi les deux extrémités.}

Notons que la \uline{catégorie opposée} à un intervalle est un intervalle, \deleted{que} et la cat. opposée à un intervalle standard \deleted{\unsure{n'}} est \deleted{\unsure{auto}} l'intervalle standard qui s'en déduit en changeant tous les signes sur les intervalles de transition $[i-1,i\deleted{\unsure{-1}}]$ ($1\leq i\leq n$). On se gardera de confondre la propriété $I\simeq I^{\circ}$, avec la propriété de symétrie.

% archive sheet number: 266
